% Generated from validated Article Draft Result. Do not hand-edit; revise the structured draft instead.
\begingroup\raggedright
\subsection{TaxBreak — inference overheadを分解する}
\par\endgroup

TaxBreakはauthorsがhost-visible orchestration overheadを、framework translation、CUDA library translation、kernel launch-pathの3要素へ分解するtrace-driven methodologyとして提案した。model FLOPSだけでは見えにくいsoftware-stack側のcostを測ろうとする研究である。\autocite{src-198727200361}


\begingroup\raggedright
\subsection{Claudini — Agentが攻撃手法そのものを探索する}
\par\endgroup

ClaudiniはauthorsがAI agentによるautoresearchでLLMへのadversarial attack algorithmを探索し、white-box jailbreakingやprompt-injection evaluationの手法を進められると報告した。Agent能力がdefenseだけでなくattack research側にも使われるという論点を提示する。\autocite{src-bb83c9c55568}


\begingroup\raggedright
\subsection{CDH-Bench / MiroEval — 何を測るかを作り直す}
\par\endgroup

CDH-Benchはauthorsがvisual evidenceとcommonsenseを意図的に衝突させ、vision-language modelのvisual fidelityを評価するbenchmarkとして提案した。MiroEvalはdeep research systemについてprocessとoutcomeを評価するframeworkを提示する。対象は異なるが、単一のfinal answer scoreだけではsystem behaviorを捉えきれないという問題意識が共通する。\autocite{src-4553a7bbd23e,src-5d84874f2d03}


\begingroup\raggedright
\subsection{ELT-Bench-Verified — benchmark自体の品質を監査する}
\par\endgroup

ELT-Bench-VerifiedはauthorsがAuditor-Corrector methodologyを提案し、LLM-driven root-cause analysisとhuman validationを組み合わせてbenchmark qualityを監査したと報告する。Agent能力の評価値を読む前に、taskやanswer keyなど測定器側の欠陥が結果へ与える影響を点検する必要があるという警告である。\autocite{src-a27679e5527b}


